% 背景知识

\section{背景知识：模态综合法 (CMS)}
给定一个离散化为 $n$ 个自由度的区域 $\Omega$，其刚度矩阵为 $\mathbf{K} \in \mathbb{R}^{n\times n}$，质量矩阵为 $\mathbf{M} \in \mathbb{R}^{n \times n}$。广义特征值问题
\begin{equation} \mathbf{K} \mathbf{u}_l= \lambda_l \mathbf{M} \mathbf{u}_l
  \label{eq:generalize-eigen-problem}
\end{equation}
给出特征值 $\lambda_l \in \mathbb{R}$ 和特征向量 $\mathbf{u}_l\in\mathbb{R}^{n}$ 的配对。模态分析通常只求解低频模态（即对应于较小 $\lambda_l$ 的模态）。

\subsection{划分}
要应用模态综合法 (CMS)，需要将区域划分为子结构。区域 $\Omega$ 的一个划分 $\mathcal{P}$ 是一组互不重叠的子结构 $\Omega_i, i=1, \cdots m$，满足 $\bigcup_i\Omega_i = \Omega$（示例见~\cref{fig:domain-partition}）。子结构之间的界面定义为 $\Gamma_{ij} = \Omega_i \cap \Omega_j$。划分 $\mathcal{P}$ 的界面定义为其并集 $\partial \mathcal{P} = \bigcup_{i,j} \Gamma_{ij}$。第 $i$ 个子结构的界面定义为 $\partial \Omega_i = \bigcup_{j} \Gamma_{ij}$。注意符号 $\partial$ 不是指外部区域边界。每个子结构 $\Omega_i$ 内部的自由度记为 $u_i \in \Omega_i\setminus\partial \Omega_i$（不包含界面），所有剩余自由度，即划分界面 $\partial \mathcal{P}$ 上的自由度，收集到向量 $u_b$ 中。不失一般性，我们假设自由度已重新排序，使得 $\mathbf{u} = [u_1^T, u_2^T, \cdots, u_m^T, u_b^T]^T$。此时 $\mathbf{K}$ 具有以下分块结构：
\begin{equation}
  \left[
    \begin{array}{ccc|c}
      \mathbf{K}_{11} & & & \mathbf{K}_{1b}\\
      &\ddots & & \vdots\\
      & & \mathbf{K}_{mm} & \mathbf{K}_{mb}\\\hline
      \mathbf{K}_{b1} &\cdots
    &\mathbf{K}_{bm} &\mathbf{K}_{bb}
    \end{array}
  \right]
  \triangleq
  \begin{bmatrix}
    \mathbf{K}_{\mathcal{I}\mathcal{I}}&\mathbf{K}_{\mathcal{I}b}\\
    \mathbf{K}_{b\mathcal{I}}&\mathbf{K}_{bb}
  \end{bmatrix}.
\end{equation}
质量矩阵 $\mathbf{M}$ 具有相同的结构。


\begin{figure}[H]
  \centering
  \newcommand{\PIC}[1]{\includegraphics[height=0.28\linewidth]{domain_definition/#1}}
  \PIC{domain-def.png}
  \PIC{inner_modes.png}
  \PIC{interface_modes.png}
  \caption{\textbf{划分、子结构特征模态和界面模态示例。} 区域 $\Omega$ 被划分为四个子结构（左图）。绘制了第 1 个子结构特征模态（中图）和一个界面模态（右图）。}
  \label{fig:domain-partition}
\end{figure}


\subsection{Craig-Bampton (CB) 方法}
我们简要介绍 CB 方法~\cite{craig1968coupling}，这是 CMS 的一种代表性实现，以概述典型的 CMS 算法流程。

\paragraph{子结构特征模态}
该方法首先通过求解下式，并行计算每个子结构的 $q_\mathcal{I}$ 个（截断后的）子结构特征模态 $\mathbf{S}_\mathcal{I}$：
\begin{equation}
  \mathbf{K}_{ii} \mathbf{S}_i = \mathbf{M}_{ii} \mathbf{S}_i \Lambda_i.
  \label{eq:sub-eig}
\end{equation}
$\mathbf{K}_{ii}$ 和 $\mathbf{M}_{ii}$ 是全局矩阵 $\mathbf{K}$ 和 $\mathbf{M}$ 的子矩阵，与由 $\Omega_i$ 中单元组装而成的刚度和质量矩阵不同。上式描述了一个约束模态分析问题，即界面 $\partial \Omega_i$ 上的自由度被固定为零。

\paragraph{界面模态}
\label{sec:interface}
随后，计算界面模态。这些模态由两部分组成：界面处的单位位移激励和子结构内部的响应，即 $\mathbf{S}_b=[\mathbf{S}_{b\mathcal{I}}^T, \mathbf{S}_{bb}^T]^T$。$\mathbf{S}_{bb} = \mathbf{I}$ 作为界面处的 Dirichlet 边界条件，而
\begin{equation}
  \mathbf{S}_{b\mathcal{I}} = - \mathbf{K}_{\mathcal{I} \mathcal{I}}^{-1}
  \mathbf{K}_{\mathcal{I}b} \mathbf{S}_{bb}= - \mathbf{K}_{\mathcal{I} \mathcal{I}}^{-1}
  \mathbf{K}_{\mathcal{I}b} \mathbf{I}
  \label{eq:traditional-response}
\end{equation}
是该边界条件下的静力平衡状态。由于 $\mathbf{K}_{\mathcal{I}\mathcal{I}}$ 是分块对角矩阵，上述方程也可以对每个子结构并行求解。

\paragraph{缩减}
\label{sec:reducing}
有了所有子结构特征模态 $\mathbf{S}_i$ 和界面模态 $\mathbf{S}_b$（见~\cref{fig:domain-partition} 中的示意），原始自由度 $\mathbf{u}$ 通过由它们组装而成的缩减矩阵 $\mathbf{S}$ 缩减为 $\mathbf{z}$：
\begin{equation}
  \mathbf{u} =
  \left[
    \begin{array}{ccc|c}
  \mathbf{S}_1 &&&\\
      &\ddots && \mathbf{S}_{b\mathcal{I}}\\
      &&\mathbf{S}_m &\\\hline
      &&&\mathbf{S}_{bb}
    \end{array}
  \right]
  \begin{bmatrix}
  z_1 \\\vdots\\ z_m\\z_b
  \end{bmatrix}
  \triangleq
  \begin{bmatrix}
    \mathbf{S}_{\mathcal{I}} & \mathbf{S}_{b}
  \end{bmatrix}\mathbf{z}
  \triangleq \mathbf{S} \mathbf{z}.
  \label{eq:redmat}
\end{equation}

应用 Galerkin 三重矩阵乘积后，原问题被缩减为一个自由度更少的特征值问题：
\begin{equation}
  \mathbf{S}^T\mathbf{K} \mathbf{S} \mathbf{z}_l =\lambda_l \mathbf{S}^T \mathbf{M} \mathbf{S} \mathbf{z}_l.
  \label{eq:reduced-matrix-pencil}
\end{equation}
该问题可以高效求解。

当计算具有以下结构的缩减刚度和质量矩阵时
\begin{equation}
  \begin{bmatrix}
    \Lambda_{1} &&&&\\
    &\ddots&&&\\
    &&\Lambda_{m}&&\\
    &&&&\tilde{\mathbf{K}}_{bb}
  \end{bmatrix},
  \begin{bmatrix}
    \mathbf{I}_{1} & & & \tilde{\mathbf{M}}_{1b}\\
    &\ddots & & \vdots\\
    & & \mathbf{I}_{m} & \tilde{\mathbf{M}}_{mb}\\
    \tilde{\mathbf{M}}_{b1} &\cdots &\tilde{\mathbf{M}}_{bm} &\tilde{\mathbf{M}}_{bb}
  \end{bmatrix},
\end{equation}
存在许多大规模矩阵-矩阵乘法：
\begin{equation}
  \begin{split}
    \tilde{\mathbf{K}}_{bb} &= \mathbf{K}_{bb} - \sum_i{\textcolor{blue}{\mathbf{K}_{bi}}
      \textcolor{red}{\mathbf{K}_{ii}^{-1} \mathbf{K}_{ib}}}\\
    \tilde{\mathbf{M}}_{ib} &= \textcolor{blue}{\mathbf{M}_{ib} -
    \mathbf{M}_{ii}} \textcolor{red}{\mathbf{K}_{ii}^{-1}\mathbf{K}_{ib}}\\
    \tilde{\mathbf{M}}_{bb} &= \mathbf{M}_{bb} - \sum_i{\color{blue} (\textcolor{red}{\mathbf{K}_{bi}
\mathbf{K}_{ii}^{-1}} \mathbf{M}_{ib} + \mathbf{M}_{bi} \textcolor{red}{\mathbf{K}_{ii}^{-1}
\mathbf{K}_{ib}})}\\
&\quad\quad\quad+\sum_i \textcolor{red}
    {\mathbf{K}_{bi}\mathbf{K}_{ii}^{-1}}\textcolor{blue}{\mathbf{M}_{ii}}
\textcolor{red}{\mathbf{K}_{ii}^{-1} \mathbf{K}_{ib}}.
  \end{split}
  \label{eq:reducing}
\end{equation}
虽然 $\mathbf{K}_{ii}^{-1}\mathbf{K}_{ib}$ 这样的项已在 $\mathbf{S}_{b\mathcal{I}}$ 中预先计算，且许多乘法（彩色文本）可以并行执行，但这仍然代价高昂。实际上，在某些文献~\cite{aristidou2013dynamic}中，界面模态计算（\cref{eq:traditional-response}）和缩减过程（\cref{eq:reducing}）统称为 Schur 补，详见 \cref{sec:schur-complement}。


\subsection{Schur 补}\label{sec:schur-complement}


总体代价主要由两部分组成：
\begin{itemize}
    \item $t_{\text{sub-eig}}$：求解子结构特征值问题（\cref{eq:sub-eig}）。
    \item $t_{\text{sub-schur}}$：\cref{eq:traditional-response}
      需要对 $\mathbf{K}_{ii}$ 进行一次分解，随后进行 $n_{b,i}$ 次前代/回代，其中 $n_{b,i}$ 是 $\partial \Omega_i$ 的自由度数。接着是 \cref{eq:reducing} 中子结构级别的矩阵-矩阵乘法（彩色项）。
\end{itemize}

如图所示，$t_{\text{sub-eig}}$（实线）与 $t_{\text{sub-schur}}$（虚线）的比较显示了两者的计算代价差异。

基于此分析，我们提出一种方法，旨在：（1）利用少量有效的子结构特征模态，（2）绕过计算代价高昂且内存密集的 Schur 补。
